% =============================================================================
% Section 5 — Results and Discussion
% PINN-ONB01 Phase 1 manuscript (IJHMT submission)
% Target length: ~2,330 words
% Final baseline model: baseline_phaseDbal (FiLM, d_z=8, 3 layers x 64, 24,005 params)
% =============================================================================

\section{Results and Discussion}\label{sec:results}

The subsections below cover three topics.
Section~\ref{subsec:verif_forward} reports code verification and
the forward $\Delta T_{\mathrm{ONB}}$ comparison against five
classical correlations. Section~\ref{subsec:uq_physics} develops a
deep-ensemble UQ analysis and audits physical-trend consistency.
Section~\ref{subsec:inverse_limits} addresses the Hsu-based inverse
cavity radius, a Simpson-type confounding, and the principal
limitations of the study.

% -----------------------------------------------------------------------------
\subsection{Verification and forward accuracy}\label{subsec:verif_forward}\label{subsec:level1}\label{subsec:forward}
% -----------------------------------------------------------------------------

The numerical implementation was verified against four canonical
problems with closed-form solutions or established correlations as
references. All four tests passed (Table~\ref{tab:level1},
Fig.~\ref{fig:level1_composite}).
Test~1 reproduces 1D steady conduction through a wall with
fixed-temperature bottom and convective top; the $L^2$ error is
$0.018\%$.
Test~2 confirms PDE residual convergence over
$N_{\mathrm{coll}}\in\{200,\ldots,5\,000\}$; the trend is monotonic
with a log--log slope of $+0.009$.
Test~3 reproduces the natural-convection boundary against
$\mathrm{Nu}=0.54\,\mathrm{Ra}^{1/4}$ (McAdams; water,
$T_\infty=\SI{293}{K}$, $L=\SI{5}{mm}$); the relative error is
$0.000\%$.
Test~4 verifies autograd derivatives of $T(z)=\sin(2\pi z)$; the
errors are $0.026\%$ and $0.011\%$.

\begin{figure}[!htbp]
  \centering
  \includegraphics[width=0.85\linewidth]{figures/fig01_level1_composite}
  \caption{Level 1 verification.
    (a) Parity between PINN and the analytical 1D conduction solution
    ($L^2$ error $0.018\%$).
    (b) PDE residual convergence with collocation density (log--log slope
    $+0.009$).
    (c) Natural-convection boundary against McAdams
    $\mathrm{Nu}=0.54\,\mathrm{Ra}^{1/4}$, $0.000\%$ relative error.
    (d) Autograd derivatives of $\sin(2\pi z)$, $0.026\%$ and $0.011\%$.}
  \label{fig:level1_composite}
\end{figure}



\begin{table}[!htbp]
\centering
\caption{Level 1 code-verification tests. All four tests pass at machine
  precision and are consistent with the expected order of accuracy of the
  spectral autograd operator.}
\label{tab:level1}
\small
\begin{tabular}{lll}
\hline
Test & Reference & Relative error \\
\hline
1D steady conduction        & Analytical linear profile & 0.018\% (L$^2$) \\
PDE residual convergence    & Monotone in $N_{\mathrm{coll}}$ & slope $+0.009$ \\
Natural-convection Nu       & McAdams correlation & 0.000\% \\
Autograd $\partial,\partial^2$ & Manufactured $\sin(2\pi z)$ & 0.026\%, 0.011\% \\
\hline
\end{tabular}
\end{table}

With the implementation verified, we turn to the forward
$\Delta T_{\mathrm{ONB}}$ comparison.
Table~\ref{tab:forward} summarizes the principal forward result. Five
classical correlations~\citep{hsu1962,davis1966,bergles1964,sato1964,basu2002}
are evaluated on the FC-77-excluded subset. The subset contains $n=77$
ONB points spanning 49 surfaces and three fluids. Bergles--Rohsenow
and Basu are valid only for water at atmospheric pressure and are
evaluated on smaller subsets, whereas the PINN is evaluated on the
full 77 because its encoder is fluid-aware.

\begin{table}[!htbp]
\centering
\caption{Forward-prediction performance on the literature ONB dataset.
  Bold row: PINN model used as the production baseline throughout the paper.
  Correlation rows quote results on their natively valid subset.}
\label{tab:forward}
\small
\begin{tabular}{lccccc}
\hline
Model & RMSE [K] & MAE [K] & $R^2$ & MRE [\%] & $n$ \\
\hline
Hsu (1962) / Sato--Matsumura (1964) & 7.43  & 6.07  & $-0.98$  & 95.0  & 43 \\
Davis--Anderson (1966)              & 16.47 & 10.26 & $-8.72$  & 334.2 & 43 \\
Bergles--Rohsenow (1964)            & 7.87  & 6.37  & $-0.87$  & 98.4  & 33 \\
Basu et al.\ (2002)                 & 7.21  & 5.85  & $-0.21$  & 106.7 & 22 \\
\textbf{PINN \texttt{baseline\_phaseDbal}} (this work)
                                    & \textbf{3.42} & \textbf{2.21}
                                    & $\mathbf{+0.44}$ & \textbf{61.4} & \textbf{77} \\
\hline
\end{tabular}
\end{table}

The PINN reduces RMSE from \SI{7.21}{K} (best classical, Basu et al.)
to \SI{3.42}{K}, a $52.6\%$ reduction. MAE falls to \SI{2.21}{K}
(from \SI{5.85}{K}), the coefficient of determination becomes
positive ($R^2=+0.44$) for the first time, and the mean relative
error contracts to $61.4\%$ from values above $95\%$ for every
classical correlation. The PINN therefore tracks
$\Delta T_{\mathrm{ONB}}$ magnitude across surfaces and fluids; the
prediction is not a per-fluid constant. Two factors bias the
comparison against the PINN. First, classical correlations are
evaluated only on their natively valid subset. Second, the PINN is
evaluated on the full $n=77$ corpus, including the more challenging
refrigerants, and on held-out points only. Parity plots are shown in
Figs.~\ref{fig:correlation_parity} (classical) and
\ref{fig:pinn_parity} (PINN).

\begin{figure}[!htbp]
  \centering
  \includegraphics[width=0.85\linewidth]{figures/fig05_correlation_parity}
  \caption{Parity plots of the five classical correlations on the
    FC-77-excluded ONB dataset. Diagonal: $y=x$; dashed: $\pm 50\%$ band.}
  \label{fig:correlation_parity}
\end{figure}


\begin{figure}[!htbp]
  \centering
  \includegraphics[width=0.65\linewidth]{figures/fig06_pinn_parity}
  \caption{Parity plot of the PINN \texttt{baseline\_phaseDbal} model on
    $n=77$ ONB points (RMSE $\SI{3.42}{K}$, $R^2=+0.44$); markers
    colored by fluid.}
  \label{fig:pinn_parity}
\end{figure}



PINN error is unevenly distributed across fluids
(Table~\ref{tab:per_fluid}). The smallest RMSE is for R-134a
(\SI{1.18}{K}, $n=34$), followed by R-123 (\SI{2.35}{K}, $n=10$).
The largest RMSE is for water (\SI{4.91}{K}, $n=33$). The water
performance reflects the heterogeneity of seven fabrication routes
at atmospheric pressure, spanning polished, EDM, plasma oxides, SAM,
and biphilic surfaces. Refrigerant data come predominantly from
Jabardo et al.\ (2009), a single laboratory whose homogeneous
fabrication route is more easily captured by the encoder.

\begin{table}[!htbp]
\centering
\caption{Per-fluid performance of the PINN \texttt{baseline\_phaseDbal} model.}
\label{tab:per_fluid}
\small
\begin{tabular}{lcccc}
\hline
Fluid & RMSE [K] & MAE [K] & $R^2$ & $n$ \\
\hline
Water   & 4.91 & 3.71 & $+0.27$ & 33 \\
R-134a  & 1.18 & 0.89 & $+0.39$ & 34 \\
R-123   & 2.35 & 1.76 & $-0.01$ & 10 \\
\hline
\end{tabular}
\end{table}

A source-paper view sharpens the diagnosis.
Table~\ref{tab:per_paper} gives the corresponding breakdown. The
smallest RMSEs are for JABARDO\_2009 (\SI{1.52}{K}) and BOURDON\_2012
(\SI{1.88}{K}). The largest are for BOURDON\_2015 (\SI{6.54}{K}) and
BETZ\_2013 (\SI{6.34}{K}). Both involve unique fabrication routes
with few replicates: silane-grafted oxides and biphilic patterns.
The model performs best on well-represented families and degrades
in proportion to the per-source sample size on under-represented
ones. This pattern motivates the UQ analysis developed below.

\begin{table}[!htbp]
\centering
\caption{Per-source-paper performance of the PINN
  \texttt{baseline\_phaseDbal} model on the 77-point ONB dataset.}
\label{tab:per_paper}
\small
\begin{tabular}{lccc}
\hline
Source paper & RMSE [K] & MAE [K] & $n$ \\
\hline
BETZ\_2013      & 6.34 & 5.91 & 10 \\
BOURDON\_2012   & 1.88 & 1.30 & 6 \\
BOURDON\_2015   & 6.54 & 4.22 & 5 \\
JABARDO\_2009   & 1.52 & 1.09 & 44 \\
JONES\_2009     & 2.81 & 1.55 & 5 \\
JO\_2011        & 3.83 & 3.75 & 2 \\
PHAN\_2009      & 4.26 & 3.87 & 5 \\
\hline
\end{tabular}
\end{table}

Three architectural features act jointly to deliver this accuracy.
First, the encoder maps heterogeneous descriptors ($R_a$,
$\theta_{\mathrm{contact}}$, optional $r_c$) into a $d_z=8$ latent,
and FiLM then injects this latent into every hidden unit. Second,
the loss balance was tuned by Bayesian optimization, yielding
$w_{\mathrm{data}}=35.57$, $w_{\mathrm{BC}}=3.89$,
$w_{\mathrm{PDE}}=2.0$, $w_{\mathrm{ONB}}=0.10$,
$w_{\mathrm{mono},R_a}=2.5$, and $w_{\mathrm{mono},\theta}=1.0$; the
tuned balance prevents the sparse ONB loss from being dominated by
the PDE residual. Third, the two-phase schedule (PDE+BC warm-up
then Adam+L-BFGS) outperformed a flat baseline by $\sim$\SI{0.5}{K}
RMSE. The $24{,}005$ parameters lie below the overfitting threshold
for the available data.

Against this baseline, the classical correlations show distinct
failure modes. Davis--Anderson has the largest residual (RMSE
\SI{16.47}{K}); it requires explicit cavity-mouth radii absent from
the corpus. Hsu and Sato--Matsumura yield identical predictions
(RMSE \SI{7.43}{K}), and Bergles--Rohsenow evaluates on $n=33$ with
RMSE \SI{7.87}{K}. Basu et al.\ yields the best $R^2$ ($-0.21$) and
the lowest classical RMSE (\SI{7.21}{K}) on its $n=22$ subset. The
PINN improves on this baseline by $52.6\%$ in RMSE and extends the
valid envelope from $n=22$ to $n=77$, a $3.5\times$ increase in the
size of the evaluable test set; it is also the only model with
positive $R^2$.

Out-of-fluid generalization shows the same pattern. On the
refrigerant subset, Hsu yields \SI{6.62}{K} (R-123) and \SI{3.60}{K}
(R-134a). The PINN yields \SI{2.35}{K} and \SI{1.18}{K},
respectively. These correspond to relative improvements of $65\%$
(R-123) and $67\%$ (R-134a), both exceeding the $46\%$ improvement
on water. The encoder isolates fluid-independent geometry, while
the non-dimensional PDE absorbs fluid-property dependence. Adding a
fluid therefore enlarges the FiLM manifold without trunk retraining.

% -----------------------------------------------------------------------------
\subsection{Uncertainty and physical consistency}\label{subsec:uq_physics}\label{subsec:uq}\label{subsec:physics}
% -----------------------------------------------------------------------------

The forward errors of Section~\ref{subsec:verif_forward} are point
predictions and carry no confidence interval. To quantify prediction
uncertainty we used the deep-ensemble approach of
\citet{lakshminarayanan2017}. $K=10$ instances of the final model
were trained from distinct seeds. Each instance underwent the
analytical warm-up and experimental fine-tuning stages on the
identical 77-point dataset, differing only in the random seed for
weight initialization. The ensemble mean and epistemic standard
deviation are
\begin{equation}
  \mu(\mathbf{x}) = \frac{1}{K}\sum_{k=1}^{K} f_k(\mathbf{x}),
  \qquad
  \sigma_{\mathrm{epi}}(\mathbf{x}) = \sqrt{\frac{1}{K-1}\sum_{k=1}^{K}
    \bigl[f_k(\mathbf{x})-\mu(\mathbf{x})\bigr]^2}.
\end{equation}
The aleatoric component is modeled as
$\sigma_{\mathrm{ale}}(\mathbf{x}) = 0.20\,|\mu(\mathbf{x})|$. This
choice is consistent with the $\pm 20\%$ run-to-run scatter reported
for ONB measurements~\citep{basu2002,jones2009}. Under independence,
$\sigma_{\mathrm{total}}^2 = \sigma_{\mathrm{epi}}^2 +
\sigma_{\mathrm{ale}}^2$, and the $95\%$ interval is
$\mu\pm 1.96\,\sigma_{\mathrm{total}}$.

Empirical 95\%-interval coverage is $98.7\%$, slightly above
nominal. The mild over-coverage reflects the fixed $\pm 20\%$
aleatoric component combined with the $K=10$ epistemic spread. The
empirical coverage exceeds the nominal $95\%$, indicating a small
but conservative bias in the intervals that is acceptable for the
engineering use of the model. An earlier \texttt{phaseB} ensemble
under the same assumption showed under-coverage ($58\%$). The
improvement here comes from the monotonicity-regularized
\texttt{phaseDbal}, which produces more diverse members in sparse
regions. The parity plot is Fig.~\ref{fig:ensemble_parity}.

\begin{figure}[!htbp]
  \centering
  \includegraphics[width=0.75\linewidth]{figures/fig07_ensemble_errorbars}
  \caption{Deep-ensemble ($K=10$) parity plot with $95\%$ intervals on
    \texttt{baseline\_phaseDbal}; empirical coverage $98.7\%$.}
  \label{fig:ensemble_parity}
\end{figure}



Decomposing the total uncertainty by component, the mean
contributions are $\bar{\sigma}_{\mathrm{epi}}=\SI{3.85}{K}$,
$\bar{\sigma}_{\mathrm{ale}}=\SI{1.34}{K}$, and
$\bar{\sigma}_{\mathrm{total}}=\SI{4.09}{K}$. Epistemic uncertainty
exceeds the aleatoric component by a factor of $\sim 2.9$.
Additional measurements, rather than tighter optimization, are
therefore the most efficient route to reducing total uncertainty.
Per-paper analysis (Fig.~\ref{fig:ensemble_std_by_category})
identifies BOURDON\_2015
($\bar{\sigma}_{\mathrm{epi}}=\SI{6.79}{K}$), BETZ\_2013
(\SI{6.75}{K}), and PHAN\_2009 (\SI{6.67}{K}) as priorities.
JABARDO\_2009 shows the smallest spread (\SI{1.96}{K}). The maximum
$\sigma_{\mathrm{total}}$ is \SI{7.92}{K} at a JONES\_2009 water
surface. Fig.~\ref{fig:ensemble_std_vs_qflux} shows no systematic
$\sigma_{\mathrm{epi}}$ broadening outside the training $q''$
envelope.

\begin{figure}[!htbp]
  \centering
  \includegraphics[width=0.85\linewidth]{figures/fig09_std_by_category}
  \caption{Epistemic uncertainty by source paper for the \texttt{phaseDbal}
    ensemble; BOURDON\_2015, BETZ\_2013, and PHAN\_2009 show the largest
    spread, JABARDO\_2009 the smallest.}
  \label{fig:ensemble_std_by_category}
\end{figure}

\begin{figure}[!htbp]
  \centering
  \includegraphics[width=0.7\linewidth]{figures/fig08_std_vs_qflux}
  \caption{Epistemic uncertainty $\sigma_{\mathrm{epi}}$ versus heat flux;
    no systematic broadening outside the training envelope.}
  \label{fig:ensemble_std_vs_qflux}
\end{figure}

Three active-learning priorities emerge from this decomposition. The
first is additional water measurements on biphilic patterns (BETZ
family). The second is within-laboratory $q''$ sweeps on single
surfaces to test $\sqrt{q''}$ per surface. The third is two or three
measurements above $200\,\si{kW/m^2}$ for high-flux electronics
cooling. The ensemble captures weight-initialization uncertainty but
not structural uncertainty (FiLM vs.\ concat, $d_z$). A model-class
ensemble is future work. A $K=3$ pilot suggests that the
\texttt{phaseDbal} epistemic spread is $10$--$15\%$ below the
\texttt{phaseB} value.

Alongside calibrated uncertainty, the model was evaluated against a
Level~3 battery of physical-consistency tests
(Table~\ref{tab:physics}). The battery covers five canonical
$\Delta T_{\mathrm{ONB}}$ trends, energy conservation, range,
extreme-input robustness, and category consistency. Eight of nine
tests passed. The mean $|\rho|$ over the trend tests is $0.90$,
indicating that the dominant physical trends have been learned
correctly.

\begin{table}[!htbp]
\centering
\caption{Level 3 physical-consistency tests for the PINN
  \texttt{baseline\_phaseDbal} model. Spearman $\rho$ is reported for the
  trend tests; thresholds are quoted in the underlying analysis script.}
\label{tab:physics}
\small
\begin{tabular}{lcc}
\hline
Test & Spearman $\rho$ / metric & Status \\
\hline
$q''\uparrow \Rightarrow \Delta T_{\mathrm{ONB}}\uparrow$ (Hsu $\sqrt{q''}$)
   & $+0.970$ & PASS \\
$R_a\uparrow \Rightarrow \Delta T_{\mathrm{ONB}}\downarrow$
   & $-0.815$ & PASS \\
$\theta_{\mathrm{contact}}\uparrow \Rightarrow \Delta T_{\mathrm{ONB}}\downarrow$
   & $-0.720$ & PASS \\
$P\uparrow \Rightarrow \Delta T_{\mathrm{ONB}}\downarrow$
   & $-1.000$ & PASS \\
$\Delta T_{\mathrm{sub}}\uparrow \Rightarrow \Delta T_{\mathrm{ONB}}\uparrow$
   & $+1.000$ & PASS \\
PDE residual (energy conservation) & mean $\approx 0.36$ & FAIL (above $10^{-3}$) \\
$\Delta T_{\mathrm{ONB}}$ range $\in [\SI{2}{K},\SI{30}{K}]$
   & 8\% violation & marginal \\
Extreme-condition robustness (NaN/Inf) & 0/4 & FAIL \\
Category consistency (CV) & 0.16 & PASS \\
\hline
\end{tabular}
\end{table}

The Hsu $\sqrt{q''}$ dependence is reproduced with $\rho=+0.970$
(Fig.~\ref{fig:physics_trends_composite}a). Earlier checkpoints
showed a sign-flip artifact, which was traced to the ONB/PDE weight
ratio and corrected in \texttt{baseline\_phaseDbal}.
$\Delta T_{\mathrm{ONB}}$ decreases monotonically with $R_a$
($\rho=-0.815$) and $\theta$ ($\rho=-0.720$); larger $\theta$
improves bubble retention. The literature angles span
$\SI{0}{\degree}$ to \SI{158}{\degree}. $R_a$ and $\theta$ are weakly
anti-correlated across sources (Pearson $r\approx-0.18$), but the
encoder receives both descriptors as independent inputs. The
monotonicity penalties enforce the physical sign on decorrelated
sweeps. Pressure and subcooling sweeps return $\rho=-1.000$ and
$\rho=+1.000$ (Fig.~\ref{fig:physics_trends_composite}d, e).

\begin{figure}[!htbp]
  \centering
  \includegraphics[width=0.95\linewidth]{figures/fig08_physics_trends_composite}
  \caption{Physics-trend audit on synthetic sweeps for \texttt{phaseDbal}.
    Spearman correlations:
    (a) $q''\uparrow\Rightarrow\Delta T_{\mathrm{ONB}}\uparrow$
    ($\sqrt{q''}$), $\rho=+0.970$;
    (b) $R_a\uparrow\Rightarrow\Delta T_{\mathrm{ONB}}\downarrow$,
    $\rho=-0.815$;
    (c) $\theta\uparrow\Rightarrow\Delta T_{\mathrm{ONB}}\downarrow$,
    $\rho=-0.720$;
    (d) $P\uparrow\Rightarrow\Delta T_{\mathrm{ONB}}\downarrow$,
    $\rho=-1.000$;
    (e) $\Delta T_{\mathrm{sub}}\uparrow\Rightarrow\Delta T_{\mathrm{ONB}}\uparrow$,
    $\rho=+1.000$;
    (f) histogram against $[\SI{2}{K},\SI{30}{K}]$.}
  \label{fig:physics_trends_composite}
\end{figure}



The energy and range tests qualify these trend results. The mean PDE
residual on a dense interior grid is approximately $0.36$ in
non-dimensional units. This value is above the $10^{-3}$ threshold
but small in absolute terms; the training-time term reached
$\mathcal{O}(10^{-8})$. The gap reflects the joint data/residual
optimization. The range test is violated for $8\%$ of synthetic
sweep points. The violations occur predominantly at extreme
low-$q''$ or high-pressure combinations outside training. The
extreme-condition test produced no NaN/Inf but one output sign
change. This is flagged as a residual failure mode outside the
trained window.

% -----------------------------------------------------------------------------
\subsection{Inverse problem and limitations}\label{subsec:inverse_limits}\label{subsec:inverse}\label{subsec:limitations}
% -----------------------------------------------------------------------------

The framework also supports inverse use: given an ONB observation,
it can recover a per-surface cavity radius $r_c$. The main inverse
result is the Hsu-based analytical inverse; a PINN-augmented variant
is reported as supplementary. For each surface and ONB observation,
the Hsu criterion was inverted to yield the range
$[r_{c,\min},r_{c,\max}]$ consistent with the measured
$\Delta T_{\mathrm{ONB}}$, saturation pressure, and subcooling. The
per-surface value is the geometric mean. Across 48 surfaces (77
observations), recovered radii span $0.39$--$\SI{13.22}{\micro\meter}$.
The mean is $\SI{3.21}{\micro\meter}$ and the median is
$\SI{2.38}{\micro\meter}$, with IQR
$[0.73,4.06]\,\si{\micro\meter}$. $60\%$ of surfaces fall within the
canonical $1$--$\SI{100}{\micro\meter}$ band. The remaining $40\%$
are predominantly sub-micron, clustering in the
refrigerant-on-stainless-steel and refrigerant-on-brass
subcategories of JABARDO\_2009.

Pooled across 48 surfaces, the $R_a$--$r_c$ correlation is weakly
negative ($\rho=-0.274$). Read at face value, this suggests that
smoother surfaces host larger cavities
(Fig.~\ref{fig:inverse_r_c_vs_Ra}).

\begin{figure}[!htbp]
  \centering
  \includegraphics[width=0.75\linewidth]{figures/fig18_inverse_rc_vs_Ra}
  \caption{Inferred $r_c$ (Hsu inverse) against $R_a$; markers colored by
    fluid. The pooled negative trend reverses sign within fluid-stratified
    subsets (Simpson's paradox).}
  \label{fig:inverse_r_c_vs_Ra}
\end{figure}

Stratifying by fluid resolves this apparent contradiction. Water
yields $\rho=+0.30$, while R-134a yields $\rho=-0.77$. The reversal
is a textbook Simpson's paradox~\citep{simpson1951} driven by
fluid-by-$q''$ confounding. Water data come from low-$q''$,
near-atmospheric operation on rougher, lower-density surfaces.
R-134a data come from high-pressure, high-$q''$ operation on
smoother, higher-density surfaces. Stratified mean radii are
$\SI{4.70}{\micro\meter}$ (water), $\SI{0.68}{\micro\meter}$
(R-134a), and $\SI{2.44}{\micro\meter}$ (R-123). These values are
consistent with higher-pressure refrigerants activating smaller
cavities. The sub-category breakdown
(Fig.~\ref{fig:inverse_r_c_by_category}) gives BOURDON\_2012
$\SI{11.73}{\micro\meter}$, JONES $\SI{5.23}{\micro\meter}$, BETZ
$\SI{1.07}{\micro\meter}$, and JABARDO\_Br $\SI{0.59}{\micro\meter}$.
These values are consistent with the order of magnitude reported for
the corresponding SEM features.

\begin{figure}[!htbp]
  \centering
  \includegraphics[width=0.85\linewidth]{figures/fig20_inverse_rc_by_category}
  \caption{Per-category distribution of $r_c$. Refrigerant categories
    cluster sub-micron; water categories span
    $1$--$\SI{15}{\micro\meter}$.}
  \label{fig:inverse_r_c_by_category}
\end{figure}



A PINN-augmented variant (\texttt{phaseDbal\_rc}) using the
Hsu-inverted $r_c$ as a weak label was trained to re-invert. For
refrigerants (R-134a 13/15, R-123 5/5) the inverse moved from the
initial guess, correlating weakly but positively with the Hsu
inverse (Spearman $\rho=+0.34$, $p=0.17$, $n=18$). For water
(30/48) the inverse did not move from initialization, indicating
insensitivity of the forward model to $r_c$ in the
water-atmospheric regime. The Hsu analytical inverse therefore
remains the headline result. The PINN-augmented variant represents
a preliminary step toward future work that will use experimentally
measured cavity distributions as hard labels.

We close by noting four limitations of the present study.

\begin{enumerate}
  \item \textit{Per-surface ONB density.} The corpus averages $1.2$ ONB
    observations per surface, precluding direct verification of
    $\sqrt{q''}$ per surface.
  \item \textit{Fluid sparsity.} The distribution is uneven (water $n=33$,
    R-134a $n=34$, R-123 $n=10$); FC-77 was excluded because CoolProp
    does not expose the fluid.
  \item \textit{No SEM/AFM ground truth.} The inverse $r_c$ cannot be
    validated against measured distributions. Reported values are
    Hsu-self-consistency estimates only.
  \item \textit{PDE-residual margin.} The energy-conservation
    residual is small in absolute terms ($\sim 0.36$ on a dense grid;
    training-time term $\mathcal{O}(10^{-8})$). It exceeds the
    $10^{-3}$ threshold. The margin is acceptable for regression, but
    tighter floors are needed for predictive control.
\end{enumerate}

These limitations map onto a coordinated experimental program of
fluid-balanced ONB sweeps, co-located scanning electron microscopy
(SEM) of the boiling surfaces, and pressure-resolved measurements on
a single fabrication batch; the active-learning priority list set
out earlier in Section~\ref{subsec:uq_physics} identifies the
categories where measurements would deliver the greatest epistemic
reduction.


% Note: PINN-augmented inverse parity figure (fig:inverse_pinn_vs_hsu) is
% relegated to the supplementary material because the corresponding result
% is reported in-text only; no inline \cref{} reference is required.
